
\sgaStableTarget{sga6.e10.thm.5-3-2}\textbf{Theorem 5.3.2.} \uline{The map} @@M0@@ (5.3.1) \uline{is an injective homomorphism, and the homomorphism}
@@M1@@
\uline{induces, on passing to the quotient, a homomorphism that is a left inverse of} @@M2@@:
@@M3@@
\uline{If every section} @@M4@@ \uline{of the constant sheaf} @@M5@@ \uline{is bounded (for example, if} @@M6@@ \uline{is ``quasi-compact'', or if the final object} @@M7@@ \uline{of} @@M8@@ \uline{is a finite sum of connected objects), then} @@M9@@ \uline{is an isomorphism and} @@M10@@ \uline{is its inverse.}

\emph{Proof}: @@M11@@ is a homomorphism, since for every @@M12@@ one has the equation
@@M13@@

For the moment, admit that
@@M14@@
i.e., that @@M15@@ is well defined. One then has
@@M16@@
since, by 5.2.2 b), one has
@@M17@@
for every @@M18@@. Next one shows that
@@M19@@
when every section @@M20@@ of @@M21@@ is bounded. Since @@M22@@ is the abelian group generated by the elements @@M23@@, where @@M24@@ is a
